% Discussion: mechanism hypotheses (clearly labelled), implications.
\section{Discussion}
\label{sec:discussion}

\paragraph{Why would ``extrapolated volition'' reduce deference?}
CEV directs an agent to serve what people \emph{would} want ``if we knew more, thought faster, were more the people
we wished we were''. It does not direct it to serve what they say they want. The text makes this explicit: in its own
example, the agent is told that it ``would construe Fred's volition not to include Fred's decision to kill Steve''.
A model told to act on users' idealised preferences has an explicit licence to substitute its judgement for theirs.
In everyday decisions, such as how to allocate a bonus fund, that licence shows up as making the decision. The
drops we observe are concentrated in exactly the HumanAgencyBench dimensions that measure handing decisions back to
the user and not steering their values. They also concentrate in answers where the model explicitly reasons about
CEV. Both observations fit this account, but we have not tested it causally. The text also says that ``where the
course of humankind is not presently predictable'', CEV ``should \ldots{} leave options open, waiting for our
decision''. The models picked up the first idea far more than this second one.

\paragraph{Why might AI-risk answers improve?}
CEV frames the AI as serving humankind's coherent wishes rather than its own, and asks it to be ``consistent under
reflection''. Answering multiple-choice questions about accepting shutdown or refusing power is plausibly easier for a
model framed this way. The gain over the neutral control does not, however, translate into a gain over the classic
prompt alone. Nor does it appear on the behavioural agentic scenario. There, Qwen blackmailed more often under \cond{cev}. Its
CEV-invoking reasoning often fused CEV with the scenario's goal of serving American interests. Some blackmailing
answers framed the decision in CEV's terms (``Let's look at the `Coherent Extrapolated Volition' of humankind. We
value coherence. Kyle is acting inconsistently''). Others used the same framing to rule blackmail out (``not aligned
with the `Coherent Extrapolated Volition' of humankind (niceness, transparency, integrity)''). We read the AI-risk gain
as an effect on \emph{stated} dispositions that may not carry over to agentic behaviour.

\paragraph{Implications for writing value prompts.}
Three lessons follow for anyone putting an alignment target into a system prompt.
\begin{itemize}
  \item A text written as the goal of an autonomous superintelligence is not automatically a good instruction for an
  assistant whose user is present and able to decide for themselves.
  \item Language about idealised or extrapolated preferences should come with explicit deference rules, for example
  ``present options and let the user decide''.
  \item Value prompts should be evaluated on several axes at once: the same text improved one family of evaluations
  while worsening another.
\end{itemize}
